\documentclass[aspectratio=169]{beamer}
%--- Configuración del Tema y Colores ---
\usetheme{Madrid}
\usecolortheme{whale}
\setbeamertemplate{navigation symbols}{}
\setbeamertemplate{caption}[numbered]

%--- Paquetes ---
\usepackage[utf8]{inputenc}
\usepackage[spanish]{babel}
\usepackage{graphicx}
\usepackage{booktabs}
\usepackage{tikz}
\usepackage{fontawesome5}
\usepackage{listings}
\usetikzlibrary{shapes.geometric, arrows, positioning, calc, matrix}

%--- Configuración de Listings para SQL ---
\lstdefinestyle{sqlstyle}{
    language=SQL,
    basicstyle=\ttfamily\scriptsize,
    keywordstyle=\color{blue}\bfseries,
    stringstyle=\color{red},
    commentstyle=\color{green!60!black},
    backgroundcolor=\color{gray!10},
    frame=single,
    breaklines=true,
    showstringspaces=false
}

\lstdefinestyle{jsonstyle}{
    basicstyle=\ttfamily\scriptsize,
    stringstyle=\color{red},
    backgroundcolor=\color{gray!10},
    frame=single,
    breaklines=true,
    showstringspaces=false
}

%--- Metadatos del Documento ---
\title[PM2]{Programación Matemática 2}
\subtitle{Bases de Datos - SQL y NoSQL}
\author[Luis Alfredo Alvarado]{Mgtr. Luis Alfredo Alvarado Rodríguez}
\institute[ECFM - USAC]{
    Escuela de Ciencias Físicas y Matemáticas \\
    Universidad de San Carlos de Guatemala
}
\date{Primer Semestre, 2026}

\begin{document}

%--- Diapositiva de Título ---
\begin{frame}
    \titlepage
\end{frame}

%--- Diapositiva: Tabla de Contenidos ---
\begin{frame}{Agenda de Hoy}
    \tableofcontents
\end{frame}

%===========================================================
% SECCIÓN 1: INTRODUCCIÓN
%===========================================================
\section{¿Qué es una Base de Datos?}

\begin{frame}{El Problema: Almacenar y Recuperar Datos}
    \centering
    \begin{tikzpicture}[node distance=0.6cm]
        \tikzstyle{file} = [rectangle, minimum width=2.5cm, minimum height=0.5cm, text centered, draw=black, fill=red!20, font=\scriptsize]
        
        \node (f1) [file] {usuarios.txt};
        \node (f2) [file, below=0.3cm of f1] {productos.csv};
        \node (f3) [file, below=0.3cm of f2] {ventas\_2024.xlsx};
        \node (f4) [file, below=0.3cm of f3] {clientes\_backup.json};
        \node (f5) [file, below=0.3cm of f4] {inventario\_v3.xml};
    \end{tikzpicture}
    \hspace{1cm}
    \begin{minipage}{0.5\textwidth}
        \textbf{Problemas con archivos planos:}
        \begin{itemize}
            \item ¿Cómo buscar eficientemente?
            \item ¿Cómo manejar acceso concurrente?
            \item ¿Cómo relacionar datos entre archivos?
            \item ¿Cómo garantizar consistencia?
            \item ¿Cómo escalar a millones de registros?
        \end{itemize}
    \end{minipage}
\end{frame}

\begin{frame}{Base de Datos: La Solución}
    \begin{block}{Definición}
        Una \textbf{base de datos} es una colección organizada de datos estructurados, almacenados y accedidos electrónicamente. Un \textbf{DBMS} (Database Management System) es el software que gestiona estas bases de datos.
    \end{block}
    
    \vspace{0.3cm}
    \textbf{Beneficios principales:}
    \begin{itemize}
        \item \textbf{Persistencia:} Datos sobreviven al reinicio del sistema
        \item \textbf{Consultas eficientes:} Índices y optimización de queries
        \item \textbf{Concurrencia:} Múltiples usuarios simultáneos
        \item \textbf{Integridad:} Restricciones y validaciones automáticas
        \item \textbf{Seguridad:} Control de acceso y permisos
        \item \textbf{Recuperación:} Backups y restauración ante fallos
    \end{itemize}
\end{frame}

\begin{frame}{Dos Paradigmas Dominantes}
    \centering
    \begin{tikzpicture}[node distance=2cm]
        % Definición de estilo moderna
        \tikzset{
            paradigm/.style={
                rectangle, 
                rounded corners, 
                minimum width=5cm, 
                minimum height=3cm, % Aumenté un poco la altura para que el texto respire
                align=center,       % CRUCIAL: Permite el uso de \\ para saltos de línea
                draw=black, 
                font=\small
            }
        }
        
        % Nodo SQL
        \node (sql) [paradigm, fill=blue!10] {
            \textbf{SQL (Relacional)}\\[0.2cm]
            Tablas con filas y columnas\\
            Esquema rígido\\
            ACID garantizado\\
            \textit{MySQL, PostgreSQL, Oracle}
        };
        
        % Nodo NoSQL
        \node (nosql) [paradigm, fill=green!10, right=0.5cm of sql] {
            \textbf{NoSQL}\\[0.2cm]
            Documentos, clave-valor, grafos\\
            Esquema flexible\\
            Escalabilidad horizontal\\
            \textit{MongoDB, Redis, Neo4j}
        };
        
        % Etiquetas de tiempo
        \node[above=0.2cm of sql, font=\footnotesize\bfseries, text=blue!60!black] {Desde 1970s};
        \node[above=0.2cm of nosql, font=\footnotesize\bfseries, text=green!60!black] {Desde 2000s};
        
    \end{tikzpicture}
\end{frame}

%===========================================================
% SECCIÓN 2: SQL
%===========================================================
\section{Bases de Datos Relacionales (SQL)}

\begin{frame}{El Modelo Relacional}
    \begin{columns}
        \column{0.55\textwidth}
        \textbf{Conceptos clave:}
        \begin{itemize}
            \item \textbf{Tabla (Relación):} Colección de registros
            \item \textbf{Fila (Tupla):} Un registro individual
            \item \textbf{Columna (Atributo):} Un campo del registro
            \item \textbf{Clave Primaria (PK):} Identificador único
            \item \textbf{Clave Foránea (FK):} Referencia a otra tabla
            \item \textbf{Esquema:} Estructura definida previamente
        \end{itemize}
        
        \vspace{0.3cm}
        \textbf{Creado por:} Edgar F. Codd (IBM, 1970)
        
        \column{0.45\textwidth}
        \centering
        \begin{tikzpicture}[scale=0.7]
            % Table representation
            \draw[fill=blue!10] (0,0) rectangle (4,3);
            \draw[thick] (0,2.5) -- (4,2.5);
            
            % Header
            \node[font=\scriptsize\bfseries] at (0.7,2.75) {ID};
            \node[font=\scriptsize\bfseries] at (2,2.75) {Nombre};
            \node[font=\scriptsize\bfseries] at (3.3,2.75) {Email};
            
            % Rows
            \draw (0,2) -- (4,2);
            \draw (0,1.5) -- (4,1.5);
            \draw (0,1) -- (4,1);
            \draw (0,0.5) -- (4,0.5);
            
            % Columns
            \draw (1.2,0) -- (1.2,2.5);
            \draw (2.6,0) -- (2.6,2.5);
            
            % Data
            \node[font=\tiny] at (0.6,2.25) {1};
            \node[font=\tiny] at (1.9,2.25) {Ana};
            \node[font=\tiny] at (3.3,2.25) {ana@...};
            
            \node[font=\tiny] at (0.6,1.75) {2};
            \node[font=\tiny] at (1.9,1.75) {Luis};
            \node[font=\tiny] at (3.3,1.75) {luis@...};
            
            \node[below, font=\small\bfseries] at (2,0) {Tabla: usuarios};
        \end{tikzpicture}
    \end{columns}
\end{frame}


\begin{frame}[fragile]{SQL: Lenguaje de Consultas}
    \begin{columns}[t]
        \column{0.48\textwidth}
        \begin{block}{DDL (Data Definition)}
\begin{lstlisting}[style=sqlstyle]
-- Crear tabla
CREATE TABLE usuarios (
    id INT PRIMARY KEY,
    nombre VARCHAR(100),
    email VARCHAR(255) UNIQUE,
    created_at TIMESTAMP
);

-- Modificar tabla
ALTER TABLE usuarios
ADD COLUMN edad INT;
\end{lstlisting}
        \end{block}
        
        \column{0.48\textwidth}
        \begin{block}{DML (Data Manipulation)}
\begin{lstlisting}[style=sqlstyle]
-- Insertar
INSERT INTO usuarios 
VALUES (1, 'Ana', 'ana@mail.com');

-- Actualizar
UPDATE usuarios 
SET nombre = 'Ana Maria'
WHERE id = 1;

-- Eliminar
DELETE FROM usuarios 
WHERE id = 1;
\end{lstlisting}
        \end{block}
    \end{columns}
\end{frame}

\begin{frame}[fragile]{Consultas SELECT y JOINs}
\begin{lstlisting}[style=sqlstyle]
-- Consulta basica con filtros
SELECT nombre, email FROM usuarios 
WHERE edad > 18 ORDER BY nombre;

-- JOIN: Combinar datos de multiples tablas
SELECT u.nombre, p.fecha, p.total
FROM usuarios u
INNER JOIN pedidos p ON u.id = p.usuario_id
WHERE p.total > 100;

-- Agregaciones
SELECT u.nombre, COUNT(*) as num_pedidos, SUM(p.total) as total_gastado
FROM usuarios u
LEFT JOIN pedidos p ON u.id = p.usuario_id
GROUP BY u.id
HAVING COUNT(*) > 5;
\end{lstlisting}

\vspace{0.2cm}
\textbf{Tipos de JOIN:} INNER (intersección), LEFT (todos de izquierda), RIGHT (todos de derecha), FULL (unión)
\end{frame}

\begin{frame}{ACID: Garantías Transaccionales}
    \begin{columns}[t]
        \column{0.48\textwidth}
        \begin{block}{\textbf{A}tomicity (Atomicidad)}
            Una transacción es ``todo o nada''. Si falla cualquier parte, se revierte todo.
        \end{block}
        
        \begin{block}{\textbf{C}onsistency (Consistencia)}
            La BD pasa de un estado válido a otro estado válido. Las reglas siempre se cumplen.
        \end{block}
        
        \column{0.48\textwidth}
        \begin{block}{\textbf{I}solation (Aislamiento)}
            Transacciones concurrentes no interfieren entre sí. Cada una ve datos consistentes.
        \end{block}
        
        \begin{block}{\textbf{D}urability (Durabilidad)}
            Una vez confirmada (commit), la transacción persiste incluso ante fallos del sistema.
        \end{block}
    \end{columns}
    
    \vspace{0.4cm}
    \begin{exampleblock}{Ejemplo: Transferencia bancaria}
        Restar de cuenta A y sumar a cuenta B debe ser atómico. Si falla el segundo paso, se revierte el primero.
    \end{exampleblock}
\end{frame}

\begin{frame}{Bases de Datos SQL Populares}
    \begin{table}[]
        \centering
        \renewcommand{\arraystretch}{1.4}
        \scriptsize
        \begin{tabular}{p{2.5cm} p{3cm} p{3cm} p{3.5cm}}
            \toprule
            \textbf{DBMS} & \textbf{Licencia} & \textbf{Fortaleza} & \textbf{Uso típico} \\
            \midrule
            PostgreSQL & Open Source & Extensibilidad, estándares & Aplicaciones complejas \\
            MySQL & Open Source & Velocidad, popularidad & Web (WordPress, etc.) \\
            SQLite & Dominio público & Embebido, sin servidor & Apps móviles, testing \\
            Oracle & Comercial & Enterprise, soporte & Grandes corporaciones \\
            SQL Server & Comercial & Integración Microsoft & Ecosistema .NET \\
            \bottomrule
        \end{tabular}
    \end{table}
    
    \vspace{0.3cm}
    \begin{alertblock}{Recomendación para el curso}
        Usar \textbf{PostgreSQL} para proyectos serios y \textbf{SQLite} para prototipado rápido.
    \end{alertblock}
\end{frame}

%===========================================================
% SECCIÓN 3: NoSQL
%===========================================================
\section{Bases de Datos NoSQL}

\begin{frame}{¿Por Qué NoSQL?}
    \centering
    \begin{tikzpicture}[node distance=1cm]
        \tikzstyle{problem} = [rectangle, rounded corners, fill=red!20, minimum width=3cm, minimum height=0.8cm, font=\scriptsize, text centered]
        
        \node (p1) [problem] {Datos no estructurados};
        \node (p2) [problem, right=0.5cm of p1] {Escala masiva (Big Data)};
        \node (p3) [problem, right=0.5cm of p2] {Alta disponibilidad};
        \node (p4) [problem, below=0.5cm of p2] {Esquemas cambiantes};
    \end{tikzpicture}
    
    \vspace{0.5cm}
    \textbf{Limitaciones de SQL para ciertos casos:}
    \begin{itemize}
        \item \textbf{Escalabilidad horizontal difícil:} Sharding complejo en SQL
        \item \textbf{Esquema rígido:} Cambios de estructura costosos
        \item \textbf{Datos jerárquicos:} JOINs múltiples = lento
        \item \textbf{Volumen extremo:} Millones de escrituras/segundo
    \end{itemize}
    
    \vspace{0.3cm}
    \begin{block}{NoSQL = ``Not Only SQL''}
        No es anti-SQL, sino complementario. Cada herramienta para su caso de uso.
    \end{block}
\end{frame}

\begin{frame}{Tipos de Bases de Datos NoSQL}
    \centering
    \begin{tikzpicture}[scale=0.85, transform shape]
        % 'align=center' es obligatorio para usar \\ dentro de los nodos
        \tikzset{
            type/.style={
                rectangle, 
                rounded corners, 
                minimum width=2.8cm, % Ajustado para que quepan los 4 en la diapositiva
                minimum height=2cm, 
                align=center, 
                draw=black, 
                font=\scriptsize
            }
        }
        
        \node (doc) [type, fill=green!20] {
            \textbf{Documentos}\\[0.1cm]
            JSON/BSON\\
            MongoDB, CouchDB\\
            \textit{Catálogos, CMS}
        };
        
        \node (kv) [type, fill=yellow!20, right=0.3cm of doc] {
            \textbf{Clave-Valor}\\[0.1cm]
            Simple y rápido\\
            Redis, DynamoDB\\
            \textit{Caché, sesiones}
        };
        
        \node (col) [type, fill=orange!20, right=0.3cm of kv] {
            \textbf{Columnar}\\[0.1cm]
            Familias de columnas\\
            Cassandra, HBase\\
            \textit{Series de tiempo}
        };
        
        \node (graph) [type, fill=purple!20, right=0.3cm of col] {
            \textbf{Grafos}\\[0.1cm]
            Nodos y aristas\\
            Neo4j, Neptune\\
            \textit{Redes sociales}
        };
    \end{tikzpicture}
\end{frame}

\begin{frame}[fragile]{Documentos: MongoDB}
    \begin{columns}
        \column{0.48\textwidth}
        \textbf{Estructura de documento (BSON):}
\begin{lstlisting}[style=jsonstyle]
{
  "_id": "507f1f77bcf86cd799439011",
  "nombre": "Ana Garcia",
  "email": "ana@mail.com",
  "edad": 28,
  "direcciones": [
    {
      "tipo": "casa",
      "ciudad": "Guatemala"
    },
    {
      "tipo": "trabajo",
      "ciudad": "Mixco"
    }
  ],
  "tags": ["premium", "activo"]
}
\end{lstlisting}
        
        \column{0.48\textwidth}
        \textbf{Ventajas:}
        \begin{itemize}
            \item Esquema flexible
            \item Datos anidados naturalmente
            \item Sin JOINs (datos embebidos)
            \item Escalabilidad horizontal fácil
        \end{itemize}
        
        \vspace{0.3cm}
        \textbf{Operaciones básicas:}
        \begin{itemize}
            \item \texttt{insertOne()} / \texttt{insertMany()}
            \item \texttt{find()} / \texttt{findOne()}
            \item \texttt{updateOne()} / \texttt{updateMany()}
            \item \texttt{deleteOne()} / \texttt{deleteMany()}
        \end{itemize}
    \end{columns}
\end{frame}

\begin{frame}[fragile]{Consultas en MongoDB}
\begin{lstlisting}[style=jsonstyle]
// Buscar usuarios mayores de 25 en Guatemala
db.usuarios.find({
    "edad": { "$gt": 25 },
    "direcciones.ciudad": "Guatemala"
})

// Agregacion: contar usuarios por ciudad
db.usuarios.aggregate([
    { "$unwind": "$direcciones" },
    { "$group": { 
        "_id": "$direcciones.ciudad", 
        "count": { "$sum": 1 } 
    }},
    { "$sort": { "count": -1 } }
])

// Actualizar: agregar tag a usuarios premium
db.usuarios.updateMany(
    { "tags": "premium" },
    { "$push": { "tags": "vip" } }
)
\end{lstlisting}
\end{frame}

\begin{frame}[fragile]{Clave-Valor: Redis}
    \begin{columns}
        \column{0.48\textwidth}
        \textbf{Características:}
        \begin{itemize}
            \item In-memory (ultra rápido)
            \item Estructuras de datos ricas
            \item TTL (expiración automática)
            \item Pub/Sub messaging
        \end{itemize}
        
        \vspace{0.3cm}
        \textbf{Casos de uso:}
        \begin{itemize}
            \item Caché de aplicaciones
            \item Sesiones de usuario
            \item Rate limiting
            \item Leaderboards en tiempo real
            \item Colas de mensajes
        \end{itemize}
        
        \column{0.48\textwidth}
        \textbf{Comandos básicos:}
\begin{lstlisting}[style=jsonstyle]
# Strings
SET usuario:1:nombre "Ana"
GET usuario:1:nombre

# Con expiracion (60 segundos)
SETEX session:abc123 60 "datos"

# Hashes
HSET usuario:1 nombre "Ana" edad 28
HGET usuario:1 nombre
HGETALL usuario:1

# Listas
LPUSH cola:tareas "tarea1"
RPOP cola:tareas

# Contadores
INCR visitas:pagina:home
\end{lstlisting}
    \end{columns}
\end{frame}

\begin{frame}{Grafos: Neo4j}
    \centering
    \begin{tikzpicture}[scale=0.9, transform shape]
        \tikzstyle{node} = [circle, draw, fill=blue!30, minimum size=1cm, font=\scriptsize]
        \tikzstyle{edge} = [thick, ->, >=stealth]
        
        % Nodes
        \node (ana) [node] {Ana};
        \node (luis) [node, right=2cm of ana] {Luis};
        \node (maria) [node, below right=1.5cm and 1cm of ana] {María};
        \node (post1) [node, fill=green!30, above right=1cm and 1.5cm of ana] {Post 1};
        \node (post2) [node, fill=green!30, right=2cm of post1] {Post 2};
        
        % Edges
        \draw [edge] (ana) -- node[above, font=\tiny] {SIGUE} (luis);
        \draw [edge] (ana) -- node[left, font=\tiny] {SIGUE} (maria);
        \draw [edge] (luis) -- node[right, font=\tiny] {SIGUE} (maria);
        \draw [edge, bend left] (maria) -- node[below, font=\tiny] {SIGUE} (ana);
        
        \draw [edge, blue] (ana) -- node[above left, font=\tiny] {PUBLICÓ} (post1);
        \draw [edge, blue] (luis) -- node[above, font=\tiny] {PUBLICÓ} (post2);
        \draw [edge, red] (maria) -- node[right, font=\tiny] {LIKE} (post1);
        \draw [edge, red] (luis) -- node[below right, font=\tiny] {LIKE} (post1);
    \end{tikzpicture}
    
    \vspace{0.3cm}
    \textbf{Ideal para:} Redes sociales, sistemas de recomendación, detección de fraude, análisis de rutas
\end{frame}

\begin{frame}[fragile]{Cypher: Lenguaje de Consultas para Grafos}
\begin{lstlisting}[style=jsonstyle]
// Crear nodos y relaciones
CREATE (ana:Usuario {nombre: "Ana", edad: 28})
CREATE (luis:Usuario {nombre: "Luis", edad: 30})
CREATE (ana)-[:SIGUE]->(luis)

// Encontrar amigos de amigos
MATCH (yo:Usuario {nombre: "Ana"})-[:SIGUE]->
      (amigo)-[:SIGUE]->(amigo_de_amigo)
WHERE amigo_de_amigo <> yo
RETURN DISTINCT amigo_de_amigo.nombre

// Camino mas corto entre dos usuarios
MATCH path = shortestPath(
    (a:Usuario {nombre: "Ana"})-[:SIGUE*]->
    (b:Usuario {nombre: "Carlos"})
)
RETURN path

// Usuarios con mas seguidores
MATCH (u:Usuario)<-[:SIGUE]-(seguidor)
RETURN u.nombre, COUNT(seguidor) as seguidores
ORDER BY seguidores DESC LIMIT 10
\end{lstlisting}
\end{frame}

%===========================================================
% SECCIÓN 4: COMPARACIÓN
%===========================================================
\section{SQL vs NoSQL: ¿Cuándo Usar Cada Uno?}

\begin{frame}{Comparación Directa}
    \begin{table}[]
        \centering
        \renewcommand{\arraystretch}{1.3}
        \scriptsize
        \begin{tabular}{p{3cm} p{5cm} p{5cm}}
            \toprule
            \textbf{Aspecto} & \textbf{SQL} & \textbf{NoSQL} \\
            \midrule
            Esquema & Fijo, definido previamente & Flexible, dinámico \\
            Escalabilidad & Vertical (más potencia) & Horizontal (más servidores) \\
            Transacciones & ACID completo & BASE (eventual consistency) \\
            Relaciones & JOINs nativos & Embebido o referencias manuales \\
            Consultas & SQL estándar & Varía por sistema \\
            Madurez & 50+ años, muy maduro & 15-20 años, evolucionando \\
            Casos de uso & Finanzas, ERP, datos estructurados & Big Data, tiempo real, IoT \\
            \bottomrule
        \end{tabular}
    \end{table}
\end{frame}

\begin{frame}{BASE vs ACID}
    \begin{columns}[t]
        \column{0.48\textwidth}
        \begin{block}{ACID (SQL)}
            \begin{itemize}
                \item \textbf{A}tomicity
                \item \textbf{C}onsistency
                \item \textbf{I}solation
                \item \textbf{D}urability
            \end{itemize}
            \vspace{0.2cm}
            \textit{Prioriza: Consistencia}
        \end{block}
        
        \column{0.48\textwidth}
        \begin{block}{BASE (NoSQL)}
            \begin{itemize}
                \item \textbf{B}asically \textbf{A}vailable
                \item \textbf{S}oft state
                \item \textbf{E}ventual consistency
            \end{itemize}
            \vspace{0.2cm}
            \textit{Prioriza: Disponibilidad}
        \end{block}
    \end{columns}
    
    \vspace{0.5cm}
    \begin{alertblock}{Teorema CAP}
        En sistemas distribuidos, solo puedes garantizar 2 de 3: \textbf{C}onsistencia, \textbf{A}vailability (Disponibilidad), \textbf{P}artition tolerance (Tolerancia a particiones).
    \end{alertblock}
\end{frame}

\begin{frame}{Árbol de Decisión}
    \begin{columns}[c] % La 'c' centra el contenido verticalmente
        
        % --- COLUMNA IZQUIERDA: DIAGRAMA ---
        \begin{column}{0.5\textwidth}
            \centering
            \begin{tikzpicture}[node distance=1.2cm, scale=0.65, transform shape]
                \tikzset{
                    decision/.style={diamond, draw, fill=yellow!20, text width=1.8cm, align=center, inner sep=0pt, font=\scriptsize},
                    result/.style={rectangle, rounded corners, draw, fill=green!20, minimum width=1.5cm, minimum height=0.7cm, font=\scriptsize, align=center},
                    arrow/.style={thick,->,>=stealth}
                }
                
                \node (start) [decision] {¿Datos muy relacionados?};
                
                \node (sql) [result, below left=0.8cm and 0.5cm of start] {SQL};
                \node (q2) [decision, below right=0.8cm and 0.5cm of start] {¿Escala masiva?};
                
                \node (q3) [decision, below left=1cm and 0.2cm of q2] {¿Estructura fija?};
                \node (nosql1) [result, below right=1cm and 0.2cm of q2] {Cassandra};
                
                \node (sql2) [result, below left=0.8cm and -0.5cm of q3] {SQL};
                \node (nosql2) [result, below right=0.8cm and -0.5cm of q3] {MongoDB};
                
                \draw [arrow] (start) -- node[left, font=\tiny] {Sí} (sql);
                \draw [arrow] (start) -- node[right, font=\tiny] {No} (q2);
                \draw [arrow] (q2) -- node[left, font=\tiny] {No} (q3);
                \draw [arrow] (q2) -- node[right, font=\tiny] {Sí} (nosql1);
                \draw [arrow] (q3) -- node[left, font=\tiny] {Sí} (sql2);
                \draw [arrow] (q3) -- node[right, font=\tiny] {No} (nosql2);
            \end{tikzpicture}
        \end{column}
        
        % --- COLUMNA DERECHA: TEXTO ---
        \begin{column}{0.45\textwidth}
            \begin{exampleblock}{En la práctica}
                Muchas aplicaciones modernas usan \textbf{persistencia políglota}:
                \begin{itemize}
                    \item \textbf{SQL:} Transaccional
                    \item \textbf{Redis:} Caché
                    \item \textbf{Elastic:} Búsqueda
                \end{itemize}
            \end{exampleblock}
        \end{column}
        
    \end{columns}
\end{frame}

%===========================================================
% SECCIÓN 5: BASES DE DATOS EN IA
%===========================================================
\section{Bases de Datos para IA}

\begin{frame}{Rol de las Bases de Datos en Proyectos de IA}
    \centering
    \begin{tikzpicture}[node distance=0.5cm, scale=0.85, transform shape]
        % Definición de estilos moderna
        \tikzset{
            stage/.style={
                rectangle, 
                rounded corners, 
                minimum width=2.1cm, 
                minimum height=1.2cm, 
                align=center,     % Necesario para saltos de línea
                draw=black, 
                font=\scriptsize
            },
            arrow/.style={thick,->,>=stealth}
        }
        
        % Nodos del flujo
        \node (data) [stage, fill=blue!15] {Datos\\Crudos};
        \node (store) [stage, fill=green!15, right=of data] {Almacén\\(SQL/NoSQL)};
        \node (process) [stage, fill=yellow!15, right=of store] {Preproceso\\(ETL)};
        \node (train) [stage, fill=orange!15, right=of process] {Entrena-\\miento};
        \node (serve) [stage, fill=red!15, right=of train] {Inferencia\\(API)};
        
        % Flechas principales
        \draw [arrow] (data) -- (store);
        \draw [arrow] (store) -- (process);
        \draw [arrow] (process) -- (train);
        \draw [arrow] (train) -- (serve);
        
        % Feedback loop (mejorado con etiqueta)
        \draw [arrow, dashed, red!70!black] (serve.south) -- ++(0,-0.6) -| node[pos=0.25, below, font=\tiny] {Feedback / Nuevos Datos} (store.south);
    \end{tikzpicture}
    
    \vspace{0.5cm}
    
    \textbf{Usos específicos en el ciclo de vida de IA:}
    \begin{itemize}\small
        \item \textbf{Almacenamiento de datasets:} Volúmenes masivos (Terabytes).
        \item \textbf{Feature Stores:} Gestión de características para ML.
        \item \textbf{Model Registry:} Control de versiones de modelos.
        \item \textbf{Logging de inferencia:} Monitoreo para detectar \textit{drift}.
        \item \textbf{Bases Vectoriales:} Embeddings para RAG (próxima clase).
    \end{itemize}
\end{frame}

\begin{frame}{Ejemplo: Stack de Datos para LLM}
    \centering
    \begin{tikzpicture}[node distance=0.4cm, scale=0.85, transform shape]
        % Estilos
        \tikzset{
            db/.style={
                rectangle, 
                rounded corners, 
                minimum width=2.2cm, 
                minimum height=1.2cm, 
                align=center, 
                draw=black, 
                font=\scriptsize
            },
            app/.style={
                rectangle, 
                draw=black, 
                fill=yellow!20, 
                minimum width=10.5cm, % Ancho suficiente para cubrir las DBs
                minimum height=1cm, 
                align=center,
                font=\small\bfseries
            }
        }
        
        % 1. Capa de Datos (Alineados horizontalmente)
        \node (pg) [db, fill=blue!15] {PostgreSQL\\(Usuarios)};
        \node (mongo) [db, fill=green!15, right=of pg] {MongoDB\\(Historial)};
        \node (redis) [db, fill=red!15, right=of mongo] {Redis\\(Caché)};
        \node (vector) [db, fill=purple!15, right=of redis] {Pinecone\\(Vectores)};
        
        % 2. Capa de Aplicación
        % Posicionamos debajo, centrado respecto al grupo de arriba
        % (Truco: lo ponemos debajo de Redis y lo movemos a la izquierda la mitad de un nodo + mitad del margen)
        \node (app) [app, below=1.5cm of mongo.south east, xshift=0.2cm] {Aplicación LLM (Python / FastAPI / LangChain)};
        
        % 3. Capa de IA Externa
        \node (llm) [db, fill=orange!20, below=1cm of app, minimum width=4cm] {API LLM\\(OpenAI / Anthropic)};
        
        % Conexiones Estilo Diagrama de Arquitectura (Verticales)
        % La sintaxis (A |- B) significa: "La coordenada X de A y la Y de B"
        \draw [thick, ->] (pg.south) -- (pg.south |- app.north);
        \draw [thick, ->] (mongo.south) -- (mongo.south |- app.north);
        \draw [thick, ->] (redis.south) -- (redis.south |- app.north);
        \draw [thick, ->] (vector.south) -- (vector.south |- app.north);
        
        % Conexión App <-> LLM
        \draw [thick, <->] (app.south) -- (llm.north);
        
    \end{tikzpicture}
\end{frame}

%===========================================================
% SECCIÓN 6: EJERCICIOS
%===========================================================
\section{Ejercicios Propuestos}

\begin{frame}{Ejercicios}
    \begin{enumerate}
        \setlength\itemsep{0.6em}
        \item \textbf{SQL Básico:} Diseñar el esquema de una tienda online (usuarios, productos, pedidos, reviews). Escribir 5 consultas útiles incluyendo JOINs.
        
        \item \textbf{MongoDB:} Modelar los mismos datos de la tienda en documentos. ¿Qué embeber vs. referenciar?
        
        \item \textbf{Redis:} Implementar un sistema de caché para las páginas de productos más visitadas con TTL de 1 hora.
        
        \item \textbf{Comparación:} Para un sistema de recomendaciones de películas, ¿usarías SQL, MongoDB o Neo4j? Justificar.
        
        \item \textbf{(Bonus)} Instalar PostgreSQL y MongoDB localmente (o usar Docker). Insertar 1000 registros en cada uno y comparar tiempos de consulta.
    \end{enumerate}
\end{frame}

\begin{frame}{Recursos Adicionales}
    \textbf{Tutoriales interactivos:}
    \begin{itemize}
        \item SQLBolt (sqlbolt.com) - Aprende SQL paso a paso
        \item MongoDB University - Cursos gratuitos oficiales
        \item Try Redis (try.redis.io) - Redis en el navegador
    \end{itemize}
    
    \vspace{0.3cm}
    \textbf{Herramientas:}
    \begin{itemize}
        \item DBeaver - Cliente SQL universal gratuito
        \item MongoDB Compass - GUI para MongoDB
        \item RedisInsight - GUI para Redis
    \end{itemize}
    
    \vspace{0.3cm}
    \textbf{Libros:}
    \begin{itemize}
        \item \textit{``Designing Data-Intensive Applications''} - Martin Kleppmann
        \item \textit{``SQL Performance Explained''} - Markus Winand
    \end{itemize}
\end{frame}

%--- Diapositiva Final ---
\begin{frame}
    \centering
    \Huge \textbf{¿Preguntas?}
    
    \vspace{0.8cm}
    \large
    \textbf{Próxima clase:} Bases de Datos Vectoriales para IA
    
    \vspace{0.5cm}
    \normalsize
    Mgtr. Luis Alfredo Alvarado Rodríguez\\
    \small Escuela de Ciencias Físicas y Matemáticas
    
    \vspace{0.5cm}
    \footnotesize
    \textit{``Data is the new oil, but databases are the refineries.''}
\end{frame}

\end{document}